% Pursuing Stacks (cotes 134-2 to 134-8): archive cote 134-5, site page 78 (PDF page 79), typescript page 265.
% typescript. Draft compared: ps4.tex:3952-4018.
% AI-transcribed by Claude Opus 5.5 (claude-opus-5-5, adaptive thinking, effort high), Message Batches API, batch ps2 (msgbatch_01NVwr1Yuo8qJf5rykeGBXPK), 2026-11-12; not yet checked by eye.
\scanpage{134-5/78}
\mspage{265}

the main properties of homotopy \xout{notions} and asphericity notions and their
interplay.

By the end of part III and the beginning of part IV, there has
been some floating in my mind as to w\over{h}{he}ther which among the two structures\hand{,}
namely contractibility structures or asphericity structures, should be
considered as ``the'' key structure for an understanding of the modelizing
story. There was a (justified) feeling, expressed first at the begin-%
ning of section 67, that in some sense, asphericity structures were
``more general'' than contractibility structures, which caused me for a
while to view them as the more ``basic'' ones. I would be more tempted at
present to hold the opposite view. \uline{The notion of a contractibility}
\uline{structure now appears as a kind of hinge between the two main sets of}
\uline{notions besides the test notions,} \uline{namely between} \xout{the} \uline{homotopy notions}
\uline{and asphericity notions}. On the one hand, as displayed in diagram (2),
the notion of a contractibility structure \xout{is the} appears as the most
stringent one among the four main \deleted{sets}\inserted{types} of homotopy \xout{notions} structures.
On the other hand, by (1) it can be equally viewed as being a special
case of \inserted{an} asphericity structure, and as such it \xout{\unsure{can be expressed b}} gives
rise to (and can be expressed by) either one of the following \deleted{four} \hand{three}
asphericity-flavored data on a category M (for any given basic localizer
$\underline{W}$ satisfying Loc 4)):

\margin{(21} \unsure{\xout{(21)}|(21)}
\[
\left\{
\begin{array}{l}
M_{\underline{W}} \subset \mathrm{Ob}\,M \quad , \quad \underline{W}_{M_c} \text{ or simply } \underline{W}_M \subset F\ell(M)\\
\phi^{\underline{W}}_{M} \text{ or } \phi_M : \deleted{$M \to \illegible{1}\ \mathrm{Hot}_{\underline{W}}$} \hand{$M \longrightarrow \mathrm{Hot}_{\underline{W}}$}\\
\deleted{$\pi_0 = \phi^{\hand{\underline{W}}\,\underline{W}_0}_{M} : M \longrightarrow (\mathrm{Sets})$} \deleted{\illegible{3}} \hand{).}
\end{array}
\right.
\]
\margin{\deleted{$\varphi^{\underline{W}}_{M}$}}

Restricting to the case when $\underline{W}$ is either $\underline{W}_{0}$ or $\underline{W}_{\infty}$, we get the corres-%
ponding structures on M , namely the three structures $M_0, W_0, \hand{$\pi$}_0$ of (20)
plus the \over{two}{three} extra structures :
\[
(22) \qquad M_\infty = M_{\underline{W}_\infty} \subset \mathrm{Ob}\,M \ , \quad W_\infty = (\underline{W}_\infty)_M \subset F\ell\,(M) \ , \ \hand{$\varphi_M : M \longrightarrow (\mathrm{Hot}$\deleted{$_{\underline{W}}$}$)$}
\]
which can be referred to as $\infty$-\uline{connected} \uline{objects} \uline{of} M, \xout{and} $\infty$-\uline{connected}
\margin{$\times$} \uline{arrows} of M\over{.}{,} \inserted{(and ``\uline{canonical functor}'' from M to Hot).} Putting together (4), (20)\hand{,}(22), we see that a contractibi-%
lity structure \inserted{$M_c$} gives rise to \deleted{eight other types of} \inserted{\deleted{nine} ten different} structures \inserted{(including $M_c$ itself)} and is deter-%
mined by each one o\over{n}{f} these, \xout{namely:}

\deleted{\xout{(23)}x \quad $J_c$ Int(M) , $W_c$ $F\ell(M)$ , \hand{\unsure{$\subset$}} $R_c$ \hand{$\subset$} $F\ell(M)$x$F\ell(M)$}
\deleted{$M_0$ \quad Ob M , \quad $W_0$}

\margin{78}
